% Audit des ajouts à mod_stage : types de question des listes d'évaluation, export par thématique
% du questionnaire pré-stage, check-list des prérequis. Relecture, scénarios et montée en charge
% dans un Moodle réel.
% Compilation : latexmk -pdf audit-types-de-question.tex
\documentclass[11pt,a4paper]{article}
\usepackage[french]{babel}
\usepackage[T1]{fontenc}
\usepackage{lmodern,microtype,booktabs,tabularx,xcolor,hyperref,enumitem}
\usepackage[margin=2.2cm]{geometry}
\usepackage{url}
\expandafter\def\expandafter\UrlBreaks\expandafter{\UrlBreaks\do\a\do\b\do\c\do\d\do\e\do\f\do\g\do\h\do\i\do\j\do\k\do\l\do\m\do\n\do\o\do\p\do\q\do\r\do\s\do\t\do\u\do\v\do\w\do\x\do\y\do\z}
\DeclareUrlCommand\code{\urlstyle{tt}}
\Urlmuskip=0mu plus 2mu
\newcommand{\grav}[1]{\textbf{\textcolor{#1}{\textbullet}}}
\newcommand{\maj}{\grav{orange}}
\newcommand{\mineure}{\grav{gray}}
\setlist[itemize]{itemsep=2pt,topsep=3pt}
\newcommand{\ok}{\textcolor{teal}{\textbf{corrigé}}}
\newcommand{\non}{\textcolor{gray}{\textbf{non traité}}}
\title{Audit des types de question et de l'export pré-stage de \texttt{mod\_stage}}
\author{Sébastien Lefebvre}
\date{10 novembre 2026}

\begin{document}
\maketitle

\section{Objet et méthode}

L'audit couvre trois ajouts au module :
\begin{itemize}
  \item cinq types de question dans les listes d'évaluation, en plus des choix, notes et
    commentaires : cases à cocher, oui/non, nombre, date et texte court ;
  \item l'export, par thématique, des réponses au questionnaire de demande (pré-stage) ;
  \item la check-list de la demande de convention, qui parle de prérequis et non plus d'objectifs.
\end{itemize}

Il combine deux approches.
\begin{enumerate}
  \item \textbf{Relecture indépendante} du code ajouté et du code qui l'appelle : cas limites,
    sécurité, performance, cohérence avec la sauvegarde, l'import entre cours et le transfert entre
    promotions, compatibilité avec Moodle~4.3 et PHP~8.0.
  \item \textbf{Tests dans un Moodle réel} : Moodle~4.5 (\code{MOODLE_405_STABLE}), PHP~8.3,
    PostgreSQL~16, serveur web intégré de PHP avec 16~processus.
    \begin{itemize}
      \item 68 scénarios fonctionnels joués par HTTP, chacun dans la session de son acteur ;
      \item un contrôle des invariants en base : forme de chaque réponse selon son type, doublons,
        réponses obligatoires vides ;
      \item trois vagues de montée en charge, de 16, 32 et 64 clients simultanés ;
      \item une mesure en volume, 20\,700 demandes sur une même thématique.
    \end{itemize}
\end{enumerate}

Les scripts sont dans \code{docs/audit/questions/}, les résultats bruts dans son sous-dossier
\code{logs/}.

\section{Synthèse}

Aucun problème critique. Les tests ont révélé quatre défauts que la relecture n'avait pas vus, dont
deux qui n'apparaissent qu'en volume ou à l'exécution. La relecture en a relevé deux majeurs et neuf
mineurs. Tous sont corrigés sauf trois, signalés comme tels. Après correction, les 68 scénarios
passent, et les invariants tiennent après les trois vagues : 22\,400 réponses, toutes de forme
valide, sans doublon.

\begin{center}
\small
\begin{tabularx}{\linewidth}{@{}l X l@{}}
\toprule
& \textbf{Constat} & \textbf{Suite} \\
\midrule
\maj & Une option contenant une balise ou un \og <\fg{} ne peut pas être choisie (tests) & \ok \\
\maj & Export pré-stage Excel : erreur 500 au-delà de quelques milliers de lignes (volume) & \ok \\
\maj & Changer le type d'une question déjà répondue rend ses réponses illisibles (relecture) & \ok \\
\maj & Transfert entre promotions : réponses rattachées à une question de type incompatible (relecture) & \ok \\
\mineure & Export pré-stage Excel : notes et nombres écrits en texte (tests) & \ok \\
\mineure & Boutons oui/non collés à l'intitulé (captures) & \ok \\
\mineure & Format d'export inconnu : exception de programmation (tests) & \ok \\
\mineure & Neuf points de la relecture (§\ref{sec:mineures}) & \ok{} (8), \non{} (1) \\
\mineure & Retirer une option d'une liste déjà répondue invalide ces réponses & \non \\
\bottomrule
\end{tabularx}
\end{center}

\section{Défauts révélés par les tests}

\paragraph{\maj{} Options contenant une balise ou un chevron.} Les réponses étaient lues avec
\code{PARAM_TEXT}, qui retire les balises. Une option « \texttt{<b>Urgences</b>} », ou
« \texttt{Poids < 5 kg} », arrivait tronquée et ne correspondait plus à aucune option : la réponse
était refusée, et le formulaire ne pouvait pas être envoyé. Le défaut existait déjà pour les choix
multiples. Les réponses à liste fermée ou à format strict (choix, cases, oui/non, note, nombre,
date) sont désormais lues telles quelles, puis contrôlées par liste blanche
(\code{stage_question_param_type()}). Les textes libres restent nettoyés.

\paragraph{\maj{} Export pré-stage Excel en volume.} Le classeur était construit en mémoire
(\code{MoodleExcelWorkbook}). Avec 20\,700 lignes, la page dépassait la limite de 256~Mo et
renvoyait une erreur 500 au bout de 23~s. Il est maintenant écrit en flux avec OpenSpout, la
bibliothèque qu'utilise Moodle pour ses propres exports : 3,5~s pour 20\,700 lignes, sans
mémoire supplémentaire.

\paragraph{\mineure{} Nombres en texte.} L'API d'export de Moodle (\code{dataformat}) convertit
toute valeur en texte pour neutraliser les formules : une note y devenait « 3 » en texte. Avec
l'écriture en flux, chaque cellule est typée. Les nombres sont des cellules numériques, et les
textes des cellules texte, jamais lues comme des formules, même quand elles commencent par
« = ». Un scénario vérifie l'absence de formule dans le classeur.

\paragraph{\mineure{} Autres.}
\begin{itemize}
  \item Les boutons oui/non suivaient l'intitulé sur la même ligne. Ils sont maintenant sur leur
    propre ligne, comme les notes.
  \item Un format d'export inconnu (\code{dataformat=inconnu}) levait une exception de
    programmation. Il est maintenant refusé par un message d'erreur.
\end{itemize}

\section{Défauts relevés par la relecture}

\paragraph{\maj{} Changement de type d'une question déjà répondue.} Seule la suppression
vérifiait l'existence de réponses. Un commentaire passé en oui/non, ou des cases à cocher passées
en choix unique, laissaient des réponses illisibles, y compris dans les autres thématiques qui
partagent la question. Ces réponses pouvaient même être perdues à la ressaisie. Une question
répondue ne change plus de type que vers un type qui accepte ses réponses : du choix unique aux
cases à cocher, du texte court au commentaire, de la note au nombre. Rien n'est enregistré
sinon, et la DEVE en est avertie. Deux scénarios vérifient le refus et la conversion permise.

\paragraph{\maj{} Transfert entre promotions.} Les questions étaient rapprochées sur leur seul
nom. Une réponse à des cases à cocher pouvait ainsi être rattachée à une question à choix unique
de même nom dans la promotion cible. Le rapprochement exige maintenant un type compatible.
Sinon, la réponse est comptée parmi les réponses abandonnées et annoncée dans l'aperçu du
transfert (\code{transfer.php} du banc).

\subsection{Points mineurs}\label{sec:mineures}

\begin{itemize}
  \item Export par thématique : stages lus par paquets de 1\,000, au lieu d'une liste \code{IN}
    qui dépasserait les limites de SQL Server et d'Oracle. \ok
  \item Accès à l'export : \code{mod/stage:view} exigé ; courriel des étudiants réservé à la DEVE,
    puisqu'aucune page ne le montre au responsable de thématique. \ok
  \item Date du formulaire de demande : calculée dans le calendrier grégorien, quel que soit le
    calendrier choisi par l'utilisateur. \ok
  \item Nombre : « 1e999 » (infini) refusé. La virgule reste lue comme séparateur décimal : c'est
    l'usage français, au prix d'une ambiguïté sur « 1,000 ». \non
  \item Tableau envoyé à la place d'une valeur (formulaire falsifié) : réponse invalide signalée,
    et non plus une erreur de programmation. \ok
  \item Libellés oui/non dans la langue de l'utilisateur ; l'anglais n'est forcé que pour le
    maître de stage dont la convention est en anglais. \ok
  \item Texte court du formulaire de demande : longueur bornée aussi côté navigateur. \ok
  \item Dates de l'export par thématique au format AAAA-MM-JJ, qui se trie. \ok
  \item Formats HTML et PDF de l'export : valeurs échappées. Formats tableur : formules
    neutralisées par le module si la version de Moodle ne le fait pas (4.3). \ok
\end{itemize}

\paragraph{Non traité.} Le formulaire d'une liste nettoie ses options (\code{PARAM_TEXT}). Une
option retirée ou modifiée, ou une option contenant une balise et réenregistrée, invalide les
réponses déjà données. Ce comportement est antérieur aux nouveaux types. Il mériterait le même
contrôle que les changements de type.

\section{Montée en charge}

Pour chaque stage, un client enchaîne quatre actions :
\begin{enumerate}
  \item auto-évaluation de l'étudiant (envoyée deux fois en même temps pour un stage sur dix,
    comme un double clic) ;
  \item avis de l'enseignant référent ;
  \item évaluation du maître de stage par son lien ;
  \item demande de convention avec le questionnaire pré-stage.
\end{enumerate}
Chaque questionnaire contient les huit types. Pendant toute la vague, deux comptes DEVE
téléchargent en boucle l'export général et l'export pré-stage. Le débogage est désactivé, comme en
production. Temps en secondes, médiane / 95\textsuperscript{e} centile.

\begin{center}
\small
\begin{tabular}{@{}l r r r@{}}
\toprule
\textbf{Action} & \textbf{16 clients} & \textbf{32 clients} & \textbf{64 clients} \\
\midrule
Auto-évaluation (8 questions) & 0,45 / 0,77 & 0,83 / 1,24 & 1,55 / 1,99 \\
Avis de l'enseignant & 0,42 / 0,75 & 0,83 / 1,27 & 1,62 / 2,12 \\
Maître de stage, envoi & 0,25 / 0,49 & 0,58 / 0,99 & 1,25 / 1,80 \\
Demande de convention, envoi & 0,53 / 0,85 & 0,84 / 1,32 & 1,89 / 2,45 \\
Export pré-stage CSV & 0,21 / 0,47 & 0,07 / 0,57 & 0,09 / 0,91 \\
Export général & 45 & 41 & 48 \\
\midrule
Durée de la vague (200 stages) & 29~s & 54~s & 60~s \\
Erreurs & 0 & 0 & 0 \\
\bottomrule
\end{tabular}
\end{center}

\begin{itemize}
  \item Aucune erreur, aucune réponse dupliquée. Les doubles clics n'ont produit qu'une
    auto-évaluation chacun : 700 événements pour 700 stages.
  \item Les temps croissent linéairement avec le nombre de clients, ce qui est attendu avec
    16~processus PHP.
  \item Quelques requêtes isolées attendent 52 à 58~s, le temps d'un export général. C'est un
    artefact du serveur intégré de PHP : il traite à la suite les connexions qu'un processus a
    acceptées. Un serveur de production (PHP-FPM) répartit les requêtes.
\end{itemize}

\paragraph{Volume.} Avec 20\,700 demandes sur une même thématique, chacune avec ses huit
réponses, l'export pré-stage prend 2,2~s en CSV et 3,5~s en Excel. Le calcul du tableau
représente 1,9~s et 67~requêtes. La liste des stages de la thématique reste à 0,1~s.

\section{Constats hors du périmètre}

Ces deux points sont antérieurs aux ajouts audités. Ils sont signalés sans être corrigés.

\begin{itemize}
  \item \textbf{Export général lent.} Pour 1\,400 stages, l'export Excel général prend 16~s hors
    charge, et 41 à 57~s pendant les vagues. Le temps se passe dans l'écriture cellule par
    cellule de la feuille « Stages » (11,5~s pour 20~requêtes SQL), car la bibliothèque Excel de
    Moodle construit tout le classeur en mémoire. Les ajouts audités n'y comptent que pour 4~\%
    (16,1~s contre 16,7~s, mesurés avec l'ancienne version de la page). Sur une activité de
    10\,000 stages, l'export approcherait les deux minutes. L'écriture en flux, adoptée pour
    l'export pré-stage, en serait le remède.
  \item \textbf{Liste falsifiée dans la demande de convention.} Une valeur absente d'une liste
    déroulante (par exemple la situation de l'étudiant) n'est pas refusée par le formulaire : elle
    disparaît des données et provoque une erreur d'écriture en base au lieu d'un message. Seul un
    formulaire falsifié peut la déclencher.
\end{itemize}

\section{Limites}

\begin{itemize}
  \item Les tests unitaires PHPUnit ont été écrits ou mis à jour, mais pas exécutés : l'installation
    de leurs dépendances n'a pas été permise dans l'environnement de test. Ils tourneront dans
    l'intégration continue du dépôt. Les fonctions qu'ils couvrent ont été exercées par les
    scénarios HTTP et par des scripts en ligne de commande dans le Moodle réel.
  \item Le pack de langue français officiel n'était pas disponible ; les chaînes du cœur de Moodle
    sont en anglais dans les tests et en français minimal dans les captures.
  \item La compatibilité avec Moodle~4.3 a été vérifiée par lecture seulement.
\end{itemize}

\end{document}
